\subsection{Convergence of the finite-dimensional dynamics}
\label{subsec:finite-dim}

Fix $d\geq1$ and a labeled and unlabeled sample with
$N,M\geq1$, and assume that the labeled sample contains
both classes. The feature vectors are arbitrary and fixed.
Under the assumptions below and a suitable constant step size,
we prove that the iterates identify a pseudo-labeling
configuration in finitely many iterations and converge
at an eventual geometric rate to the minimizer of the
corresponding branch objective.

\begin{assumption}[Shape and curvature of loss and penalty]
\label{assumption:finite-dim-loss-penalty}
We assume the following:
\begin{enumerate}
    \item For each $y\in\{-1,+1\}$, the function
    $r\mapsto\ell(y,r)$ belongs to
    $C^2(\mathbb R;[0,\infty))$. There exists
    $L_\ell\in(0,\infty)$ such that
    \[
        y\ell'(y,r)<0,
        \qquad
        0<\ell''(y,r)\leq L_\ell,
        \qquad
        y\in\{-1,+1\},\quad r\in\mathbb R.
    \]

    \item The penalty $j$ belongs to
    $C^2(\mathbb R;[0,\infty))$ and is coercive.
    There exists $L_j\in(0,\infty)$ such that
    \[
        0<j''(u)\leq L_j,
        \qquad u\in\mathbb R.
    \]
\end{enumerate}
\end{assumption}

\subsubsection{Potential formulation}

Write
\[
    L_{\mathcal L}(\theta)
    :=
    \frac1N\sum_{i\in I_{\mathcal L}}
        \ell(Y_i,r_i(\theta)),
    \qquad
    G(\theta)
    :=
    \nabla_1L_{\mathrm{ST}}(\theta;\theta).
\]
Here the derivative is taken with the teacher parameter held fixed.
The recursion therefore becomes
\begin{equation}
    \theta^{t+1}
    =
    \theta^t-\eta P_dG(\theta^t),
    \qquad t\geq0.
    \label{eq:ST-recursion-G}
\end{equation}
Although $G$ is discontinuous at the selection thresholds,
it agrees almost everywhere with the gradient of a continuous
potential $V$. The selection convention specifies the update
at points where this potential may fail to be differentiable.

\begin{proposition}[Self-training potential]
\label{prop:ST-potential}
Under Assumption~\ref{assumption:finite-dim-loss-penalty}, define
\[
    \psi_\kappa(r)
    :=
    \mathbf1_{\{r\leq\kappa_-\}}\ell'(-1,r)
    +\mathbf1_{\{r\geq\kappa_+\}}\ell'(+1,r),
    \qquad
    \Psi_\kappa(r)
    :=
    \int_0^r\psi_\kappa(u)\,du.
\]
Then $\Psi_\kappa$ is bounded and globally Lipschitz.
The function $V:\mathbb R^{d+1}\to\mathbb R$ defined by
\[
    V(\theta)
    :=
    L_{\mathcal L}(\theta)
    +\pi V_{\mathcal U}(\theta)
    +\frac{\lambda}{d}J(w),
    \qquad
    V_{\mathcal U}(\theta)
    :=
    \frac1M\sum_{j\in I_{\mathcal U}}
        \Psi_\kappa(r_j(\theta))
\]
is continuous and twice continuously differentiable outside
the threshold set
\[
    \mathcal S_\kappa
    :=
    \bigcup_{j\in I_{\mathcal U}}
    \left(
        \{\theta:r_j(\theta)=\kappa_-\}
        \cup
        \{\theta:r_j(\theta)=\kappa_+\}
    \right),
\]
which is a finite union of proper affine hyperplanes. For every
$\theta\notin\mathcal S_\kappa$,
\[
    \nabla V(\theta)=G(\theta).
\]
\end{proposition}

To describe the branch structure, we record the selected
pseudo-labels.
A configuration is an element of
$\{-1,0,+1\}^{I_{\mathcal U}}$.
Define the configuration map
\[
    \mathfrak C:\mathbb R^{d+1}\to\{-1,0,+1\}^{I_{\mathcal U}}, \qquad 
    \mathfrak C(\theta)
    :=
    \bigl(S_j(\theta)\hat Y_j(\theta)\bigr)_{j\in I_{\mathcal U}}.
\]
For each configuration $\mathfrak c$, define its branch by
\[
    \Theta_{\mathfrak c}
    :=
    \{\theta\in\mathbb R^{d+1}:
        \mathfrak C(\theta)=\mathfrak c\},
\]
and its smooth branch objective
$L_{\mathfrak c}:\mathbb R^{d+1}\to\mathbb R$ by
\[
    L_{\mathfrak c}(\theta)
    :=
    L_{\mathcal L}(\theta)
    +
    \frac{\pi}{M}
    \sum_{\substack{j\in I_{\mathcal U}\\\mathfrak c_j\neq0}}
        \ell(\mathfrak c_j,r_j(\theta))
    +
    \frac{\lambda}{d}J(w).
\]

\begin{proposition}[Configuration and branch formulation]
\label{prop:configuration-branch-formulation}
For each configuration $\mathfrak c$, there exists a constant
$K_{\mathfrak c}$ such that
\[
    V=L_{\mathfrak c}+K_{\mathfrak c}
        \quad\text{on}\quad\overline{\Theta_{\mathfrak c}},
    \qquad
    G=\nabla L_{\mathfrak c}
        \quad\text{on}\quad\Theta_{\mathfrak c}.
\]
Each branch is convex, possibly empty, and is described by finitely
many strict or non-strict affine inequalities.
Each nonempty branch has a polyhedral closure, and the nonempty
branches partition $\mathbb R^{d+1}$ into at most $3^M$ regions.
\end{proposition}
Consequently, the self-training recursion satisfies
\[
    \theta^{t+1}
    =
    \theta^t-\eta P_d
    \nabla L_{\mathfrak C(\theta^t)}(\theta^t),
    \qquad t\geq0,
\]
including at selection thresholds.
\begin{proposition}[Branch curvature and coercivity]
\label{prop:branch-convexity-coercivity}
Suppose Assumption~\ref{assumption:finite-dim-loss-penalty} holds.
Assume that $\lambda>0$, $\pi\geq0$, $\kappa_-<0<\kappa_+$,
and both classes occur among the labeled observations.

Then every branch objective $L_{\mathfrak c}$ has a
positive-definite Hessian on $\mathbb R^{d+1}$ and is therefore
strictly convex. Moreover, $L_{\mathfrak c}$ and the potential
$V$ are coercive:
\[
    L_{\mathfrak c}(\theta)\longrightarrow+\infty,
    \qquad
    V(\theta)\longrightarrow+\infty
    \qquad\text{as }\|\theta\|\longrightarrow\infty.
\]
Consequently, each branch objective admits a unique global minimizer,
denoted by
\[
    \theta_{\mathfrak c}^\star
    :=
    \operatorname*{argmin}_{\theta\in\mathbb R^{d+1}}
        L_{\mathfrak c}(\theta).
\]
\end{proposition}


\subsubsection{Algorithmic convergence}

To state the step-size condition, define
\begin{align}
    H_d
    &:=
    \frac{L_\ell}{N}\sum_{i\in I_{\mathcal L}}
        \breve X_i\breve X_i^\top
    +
    \frac{\pi L_\ell}{M}\sum_{j\in I_{\mathcal U}}
        \breve X_j\breve X_j^\top
    +
    \frac{\lambda L_j}{d}
    \begin{bmatrix}
        0 & 0_d^\top\\
        0_d & I_d
    \end{bmatrix}, \label{eq:hessian} \\
    \Lambda_d
    &:=
    \lambda_{\max}(P_d^{1/2}H_dP_d^{1/2})>0. \label{eq:lambda-d}
\end{align}
By Proposition~\ref{prop:branch-convexity-coercivity},
for every configuration $\mathfrak c$ and
$\theta\in\mathbb R^{d+1}$, the branch Hessians satisfy
\[
    0\prec\nabla^2L_{\mathfrak c}(\theta)\preceq H_d.
\]
We measure errors in the norm
\[
    \|\theta\|_{P_d^{-1}}
    :=
    \sqrt{\theta^\top P_d^{-1}\theta}, \quad \theta \in \mathbb{R}^{d+1}.
\]
Define
\begin{align}
    T_{\mathrm{id}}
    :=
    \inf\left\{
        T\in\mathbb N_{\geq1}:
        \mathfrak C(\theta^t)=\mathfrak C(\theta^T)
        \text{ for every integer }t\geq T
    \right\}, \label{eq:identification-time}
\end{align}
with $\inf\varnothing=+\infty$.

\begin{theorem}[Convergence of self-training]
\label{thm:st-convergence}
Fix $d\geq1$ and a finite labeled and unlabeled sample with
$N,M\geq1$. Suppose
Assumption~\ref{assumption:finite-dim-loss-penalty} holds.
Suppose also that $\lambda>0$, $\pi>0$, $\kappa_-<0<\kappa_+$,
both classes occur among the labeled observations, and
\begin{equation}
    0<\eta<\frac{2}{\Lambda_d}.
    \label{eq:st-descent-step-size}
\end{equation}
Then, for every initial state $\theta^0\in\mathbb R^{d+1}$,
the identification time \eqref{eq:identification-time} is finite: $T_{\mathrm{id}} < \infty$, and there exists
a configuration $\mathfrak c$ such that
\[
    \mathfrak C(\theta^t)=\mathfrak c
    \quad\text{for all }t\geq T_{\mathrm{id}},
    \qquad
    \theta^t\xrightarrow[t\to\infty]{}
    \theta_{\mathfrak c}^\star
    \in\overline{\Theta_{\mathfrak c}}.
\]
Moreover, there exists $q\in(0,1)$ such that
\[
    \|\theta^t-\theta_{\mathfrak c}^\star\|_{P_d^{-1}}
    \leq
    q^{\,t-T_{\mathrm{id}}}
    \|\theta^{T_{\mathrm{id}}}-\theta_{\mathfrak c}^\star\|_{P_d^{-1}},
    \qquad t\geq T_{\mathrm{id}}.
\]
\end{theorem}

\begin{remark}
Let $\mathfrak c$ be the configuration identified in
Theorem~\ref{thm:st-convergence}.
Then $G(\theta_{\mathfrak c}^\star)=0$ if and only if
$\theta_{\mathfrak c}^\star\in\Theta_{\mathfrak c}$.
In particular, this holds whenever
$\theta_{\mathfrak c}^\star\notin\mathcal S_\kappa$.
\end{remark}

The proof of Theorem~\ref{thm:st-convergence} proceeds in three steps
and is detailed in Appendix~\ref{sec:proof-finite-dim}:
\begin{enumerate}
    \item A global descent inequality for $V$, valid across
    selection thresholds, gives $G(\theta^t)\to0$, while
    coercivity confines the trajectory to a bounded sublevel set.

    \item If configurations changed infinitely often, finiteness
    of the configuration set, boundedness, and vanishing increments
    would yield a shared accumulation point at which two distinct
    branch gradients vanish. Separation of these gradients on
    shared boundaries rules this out.

    \item Every accumulation point is then a stationary point of
    the identified strictly convex branch objective. Uniqueness
    gives convergence of the full sequence, and positive
    definiteness and continuity of the branch Hessian at the
    limit give the geometric rate.
\end{enumerate}
